\documentclass[11pt]{article}

\usepackage[utf8]{inputenc}
\usepackage{amsmath, amssymb}
\usepackage{graphicx}
\usepackage{booktabs}
\usepackage{geometry}
\usepackage{hyperref}
\usepackage{natbib}
\geometry{letterpaper, margin=1in}
\bibliographystyle{apalike}

\title{The Informational Bifurcation of Hidden Liquidity:\\
Aggressive versus Midpoint Concealment on a Single Venue}
\author{Nicholas Jiang}
\date{August 2026}

\begin{document}

\maketitle

\begin{abstract}
Hidden liquidity is commonly treated as a single object, and its informativeness as a single
number. Using event-by-event LOBSTER message data for 474 NASDAQ-traded names over
2023--2024, we show that it is neither. Hidden executions away from the midpoint---trades
with an identifiable aggressor---are associated with a signed midpoint move of $+2.09$ bps
that is flat from three minutes to the closing bell across the full cross-section. Hidden
executions \emph{at} the midpoint carry a small \emph{negative} markout of $-0.42$ bps,
drifting to $-0.71$ by thirty minutes. Pooling them, as any classifier that force-signs
midpoint fills must, averages an informed population against an uninformed one and returns a
statistical zero---which is why the four standard trade-sign rules disagree about the
informativeness of hidden liquidity by two full basis points on the same data. The negative
leg is not an artifact of the tick rule used to sign it: a matched placebo applying the
identical rule to decoupled moments recovers $5.4\%$ of it and is statistically zero, and
tick-signing the aggressive population drives its footprint to zero rather than negative. The
split is the empirical signature of the self-selection equilibrium of \citet{zhu2014},
measured \emph{within} one venue and one order type, where venue choice, fee schedule and
clientele are held fixed and only execution aggressiveness varies---a sharper test of that
mechanism than the cross-venue comparisons on which the dark-pool literature relies.

The aggressive footprint behaves as an adverse-selection quantity. It is a stable fraction of
the quoted half-spread: the cross-name correlation is $+0.85$, the ratio is centred on $0.34$
with an interquartile range of $0.28$--$0.42$, and across spread quintiles spanning a factor
of five in width the share moves only between $0.29$ and $0.38$. Entered jointly with the
spread, realized volatility---which predicts the footprint strongly on its own---falls to
exactly zero ($t=0.2$), so the footprint tracks the spread rather than volatility. On the
\citet{huangstoll1997} accounting it is $51\%$ of the effective half-spread a taker pays. It
is incremental to visible order-flow imbalance \citep{cont2014}, and it survives the
exclusion of earnings windows and large-move days with $94\%$ retention.

Three things this design cannot settle, which we state rather than absorb. Displacement:
$30\%$ of aggressive prints are followed by the touch being swept within $100$ms, and that
minority carries $75\%$ of the measured level, essentially all of it coincident with the
sweep---at this resolution, mechanical displacement and information taken through the trade
are the same event, though the unswept $70\%$ still carry $+0.46$ bps. Causality: $47.5\%$ of
the signed move precedes the print on unambiguously-signed trades, and a Hasbrouck VAR with
thirty minutes of genuine memory finds no lasting response, so we report an association
rather than an impact---while noting that the obvious bridge between those facts fails, since
decomposing the markout against the pre-print drift leaves $+1.407$ of $+1.379$ bps
orthogonal to it. Venue: because LOBSTER reconstructs the NASDAQ book alone, we cannot
separate price discovery from a local midpoint adjusting toward a consolidated quote that
moved first. The bifurcation is immune to all three, being a contemporaneous comparison
between two populations observed on the same feed; it is the magnitudes that are provisional.
Along the way we document four ways this measurement misleads if taken at face value---a VAR
whose reported horizons lie beyond its own memory, a non-overlapping sort whose sign depends
on sampling phase, a sweep test whose conditioning window overlaps its outcome, and a
tick-rule placebo that turns out to clear rather than convict.
\end{abstract}

\section{Introduction}
\label{sec:intro}

That hidden orders carry information is established \citep{bessembinder2009, hautsch2012}.
That informed traders sort themselves across execution mechanisms according to their demand
for immediacy is likewise established, in theory by \citet{zhu2014} and empirically across
trading venues by \citet{comertonforde2015} and \citet{buti2017}. This paper joins the two.
It asks whether the sorting is visible \emph{inside} a single venue and a single order
type, and finds that it is---sharply enough that the informativeness of hidden liquidity
cannot be summarized by one number.

The empirical object is the type-5 print in the NASDAQ ITCH message stream: an execution
against non-displayed resting liquidity. Roughly half of these prints execute exactly at
the prevailing midpoint and half away from it. The two halves behave differently in a way
that is not a matter of degree. Prints away from the midpoint are associated with a signed
three-minute midpoint move that is flat across every horizon we measure---out to thirty
minutes, and, against a time-of-day-stratified placebo, all the way to the closing bell on
the full 470-name panel. Prints at the midpoint
carry $-0.42$ bps, drifting to $-0.71$ by thirty minutes. Because the standard trade-sign
classifiers differ chiefly in how they handle midpoint trades, they disagree on the
aggregate footprint by two full basis points---from $+2.09$ under a rule that abstains on
them to a statistical zero under the Ellis--Michaely--O'Hara hybrid \citep{ellis2000},
which force-signs them. That disagreement is not a robustness failure to be adjudicated. It is the
bifurcation showing through the classifier.

Our contribution is threefold, and the third is a negative result about our own measurement
instrument.

First, we measure the sorting within a venue. A cross-venue comparison confounds
self-selection with venue-level differences in fee structure, clientele, and reporting
latency. Here, venue, order type, and reporting mechanism are held fixed, and the only
variable is whether the execution crossed the spread. We regard this as a sharper test of
the \citet{zhu2014} mechanism than the cross-venue evidence, and we are explicit that the
mechanism itself is not our discovery.

Second, we characterize the aggressive footprint as the adverse-selection component of the
spread. The evidence is a set of cross-sectional facts rather than a structural estimate:
the footprint is a stable fraction of the quoted half-spread across a fivefold range of
spread widths; volatility, which predicts it univariately, contributes nothing once the
spread is controlled; it constitutes $51\%$ of the effective half-spread on the
\citet{huangstoll1997} accounting; it is incremental to visible order-flow imbalance; and
it survives the removal of information-event windows. A direct test of
mechanical displacement is only partly favourable, and we report it as such. We are careful to distinguish the facts from the
interpretation. A stable ratio is consistent with the Glosten--Milgrom mechanism \citep{glosten1985} under
which the spread compensates providers for informed flow; it does not demonstrate that the
spread is \emph{set} that way, and we do not test the causal direction.

Third, we show what our own headline measurement overstates. Conditioning on aggressive
same-side execution may select moments at which the midpoint is already travelling in that
direction, in which case a forward markout measures continuation rather than response. We
test this directly with an event-time profile on both sides of the print, restricted to
prints executing outside the prevailing quote so that the aggressor's side cannot be
misassigned. Nearly half the signed move---$47.5\%$ over $[-30\text{s},+180\text{s}]$---
precedes the print. A Hasbrouck VAR with genuine memory at the horizons of interest, which
conditions on the prior return path, independently finds no lasting response. We withdraw
the causal reading accordingly, and report the association. We then test the mechanism that
would explain both facts---that the markout is continuation of a move already underway---and
find that it does not hold: regressing the forward markout on the pre-print drift leaves
essentially all of it orthogonal to the prior path, with the continuation component
economically zero and slightly negative. Half the move precedes the print, but the pre-drift
does not forecast what follows, and the markout--VAR disagreement therefore remains open
rather than resolved.

Section~\ref{sec:data} describes the data and the burst construction. Section~\ref{sec:lit}
reports the negative result that motivates the rest: on two pilot names, clustering
\emph{visible} aggressive flow yields no informed footprint under any parameterization we
tried or under two alternative reconstructions. That section is pilot-scale throughout and
we mark it as such. Section~\ref{sec:bifurcation} establishes the bifurcation.
Section~\ref{sec:adverse} characterizes the aggressive footprint as an adverse-selection
quantity. Section~\ref{sec:alternatives} tests the two leading non-informational accounts.
Section~\ref{sec:overstate} reports what unconditional markouts get wrong.
Section~\ref{sec:limits} states what this design cannot settle.

\section{Data and Construction}
\label{sec:data}

\subsection{Data}

The analysis uses event-by-event limit-order-book data from LOBSTER, reconstructed directly
from the NASDAQ ITCH message stream, over 2023--2024. The panel is 474 names with
sufficient hidden flow, spanning $237{,}948$ raw ticker-days, of which roughly $40\%$ are
absent from the staged archive or fail book reconstruction. Each test then applies the
filter its estimator requires, so analysis-ready name-day counts differ across tables and
are reported in each caption. No name-day is excluded on the basis of its outcome.

Two properties of the data govern what can be claimed.

\textbf{LOBSTER does not disclose the side of hidden executions.} The Direction field is
uniformly $+1$ for type-5 prints, so hidden trades carry no native sign and must be
classified. This is not incidental to the paper; it is the reason the bifurcation is
visible at all, since it forces an explicit decision about midpoint prints that a signed
feed would obscure.

\textbf{LOBSTER reconstructs the NASDAQ book only.} All midpoints, quotes and depths are
NASDAQ's, not the consolidated NBBO. Section~\ref{sec:limits} states the consequence, which
is material.

Because each ticker-day must be streamed from the raw archive, decompressed and
reconstructed individually, we restrict to a two-year window rather than a longer one.

\subsection{Burst construction}
\label{sec:burstdef}

Isolated order submissions are noisy at horizons beyond a few seconds, so we aggregate
temporally clustered same-side executions into discrete \emph{bursts}. For burst $b$ with
start time $t_b$ and reference midpoint $m_{t_b}$, the signed aggregate volume is
$Q_b = \sum_{j \in b} \text{side}_j\, q_j$.

Cluster membership is determined by a \emph{self-exciting counter}, which we describe
precisely because its status is often overstated in this literature. A running intensity
decays exponentially between events and is incremented by each arrival:
\begin{equation}
\lambda(t_k) \;=\; \lambda(t_{k-1})\, e^{-\beta (t_k - t_{k-1})} \;+\; 1,
\qquad \lambda(t_0) = 1,
\label{eq:counter}
\end{equation}
and the burst remains active while $\lambda$ exceeds a threshold $\lambda_{\min}$. This is
the intensity recursion of an exponential-kernel Hawkes process, and we use it as a
\emph{clustering heuristic}; we do not fit it. No likelihood is maximized, no branching
ratio is estimated, and no goodness-of-fit is reported, so we do not describe the
construction as a fitted point process and none of our claims rest on its being one. What
we do verify is that the clustering it detects is not an arrival-rate artifact: the count
of same-side clustered bursts exceeds an \emph{inhomogeneous}-Poisson null matched to each
name's empirical intraday intensity profile by a median $z$ of $47$, rejecting the Poisson
null on $100\%$ of name-days.

Equation~\eqref{eq:counter} has two free parameters, not three. The excitation increment is
fixed at unity, which is a normalization rather than a choice: scaling the increment and
$\lambda_{\min}$ together leaves the detector's behaviour unchanged, so only their ratio is
identified. The decay rate $\beta$ and the threshold $\lambda_{\min}$ are the genuine
parameters, and both are varied---$\beta \in \{0.5, 1, 2, 5\}$ and
$\lambda_{\min} \in \{0.3, 0.5, 0.8\}$---in Section~\ref{sec:lit}. Additional filters on
fractional sweep volume, directional consistency, and opposing-volume ratio are applied
uniformly and described in the replication materials.

Bursts of \emph{hidden} executions, which carry the paper's results, are formed identically
on type-5 prints alone.

\subsection{Inference}

Throughout, the day---not the individual print---is the unit of statistical inference.
Markouts are equal-weighted within name-day, then averaged across the cross-section, and
$t$-statistics are Newey--West \citep{neweywest1987} with 10 lags on the resulting daily
mean series, or day-level Fama--MacBeth \citep{famamacbeth1973} with the same correction
where a regression is reported. An outlier rule is applied before any figure is quoted:
name-days whose markout exceeds $1{,}000$ bps in absolute value at any horizon are dropped.
On the primary panel this removes 56 name-days of $221{,}261$ ($0.03\%$); one reaches
$-114{,}627$ bps, which is a corrupt price rather than a return.

\section{A Two-Name Null for Clustered Visible Flow}
\label{sec:lit}

We begin with the negative result that directs the paper toward hidden liquidity, and state
its scope before its content: everything in this section is measured on two pilot names,
AAPL and TSLA. It is what motivated the reconstruction choice the rest of the paper makes,
not a cross-sectional claim, and we do not extend it to one.

Clustering \emph{visible} aggressive executions produces no measurable informed footprint,
and this does not depend on how the clustering is tuned. Table~\ref{tab:hawkes_grid}
recomputes the three-minute directional markout over the decay grid
$\beta \in \{0.5,1,2,5\}$; every cell is economically negligible and well inside the
bid--ask spread, and minimum cluster size (over $\{3,5,10\}$, omitted) has no effect. The
threshold $\lambda_{\min}$ was varied over $\{0.3,0.5,0.8\}$ during detector calibration
with the same result. The null is not an artifact of a particular parameterization.

\begin{table}[htbp]
\centering
\begin{tabular}{lcccc}
\toprule
\textbf{3-min markout (bps)} & $\beta=0.5$ & $\beta=1$ & $\beta=2$ & $\beta=5$ \\
\midrule
AAPL & $+0.02$ & $-0.09$ & $-0.03$ & $+0.14$ \\
TSLA & $-0.10$ & $+0.07$ & $-0.02$ & $+0.26$ \\
\bottomrule
\end{tabular}
\caption{Decay-rate grid for the visible-flow clustering rule, evaluated ungated. Every
configuration is sub-spread. Minimum cluster size $\in \{3,5,10\}$ is omitted as it had no
effect.}
\label{tab:hawkes_grid}
\end{table}

Because tuning does not help, we replace the construction entirely with three alternative
reconstructions of the same message stream: \emph{order-flow imbalance} (OFI; the
top-of-book queue dynamics of \citealp{cont2014}, and its deep-learning descendants,
\citealp{kolm2023}), \emph{book resilience} (aggressive sweeps whose consumed depth fails
to replenish), and \emph{hidden-execution} clustering. Table~\ref{tab:altdefs} reports the
comparison. Only the hidden-execution reconstruction produces a positive markout on both
names.

\begin{table}[htbp]
\centering
\begin{tabular}{lcccc}
\toprule
\textbf{3-min markout (bps)} & Clustered visible & OFI & Book resilience & Hidden \\
\midrule
AAPL & $-0.08$ & $+0.08$ & $+0.19$ & $+0.42$ \\
TSLA & $+0.08$ & $+0.08$ & $+0.16$ & $+0.31$ \\
\bottomrule
\end{tabular}
\caption{Four reconstructions of the same message stream, evaluated ungated on the 2023
sample. Hidden bursts are signed by the quote rule (execution price against the prevailing
midpoint); the other three use aggregated signed flow. Two names only: OFI and
book resilience were validated at pilot scale, and we make no cross-sectional claim for
them.}
\label{tab:altdefs}
\end{table}

Two construction choices are dictated by the definitions and applied uniformly. The
book-resilience non-refill condition is defined on queue \emph{depth} rather than quote
price, and its markout is measured strictly after the depth-observation window closes, so
the filter cannot select on the path over which it is subsequently evaluated. And hidden
prints, being unsigned in the feed, are signed by the quote rule throughout.

The rest of the paper concerns the one reconstruction that survives. The scope limit bears
repeating because the sections that follow are cross-sectional and this one is not: all four
reconstructions were compared on two names, so this is the reason we study hidden-execution
clustering, not a ranking of the four at cross-sectional scale. A referee is entitled to
read Section~\ref{sec:lit} as motivation rather than evidence, and nothing downstream
depends on it.

\section{The Bifurcation}
\label{sec:bifurcation}

\subsection{Classifier disagreement}

The headline hidden-execution footprint is measured with the quote rule, which is undefined
for the $48.9\%$ of type-5 prints executing exactly at the midpoint and therefore abstains
on them. Since that convention discards half the data, the natural question is whether the
footprint survives classifiers that instead \emph{assign} a sign to midpoint prints.
Table~\ref{tab:signing} re-signs under four canonical rules.

\begin{table}[htbp]
\centering
\begin{tabular}{lcc}
\toprule
\textbf{Classifier (3-min hidden markout)} & \textbf{Markout (bps)} & \textbf{NW $t$} \\
\midrule
Quote rule, abstaining at the midpoint & $+2.09$ & $+25.3$ \\
Tick rule (midpoint prints signed by last price change) & $+0.62$ & $+9.5$ \\
CLNV \citep{chakrabarty2007} & $+0.46$ & $+8.3$ \\
EMO \citep{ellis2000} & $+0.03$ & $+0.6$ \\
\bottomrule
\end{tabular}
\caption{Three-minute hidden-execution footprint under four trade-sign classifiers
(474 names, 2023--2024, $199{,}627$ name-days; $50\%$ of prints at the midpoint). The
aggregate ranges over two basis points depending solely on the treatment of midpoint
prints.}
\label{tab:signing}
\end{table}

The disagreement is not diffuse. It is concentrated entirely in names whose hidden prints
are predominantly at the midpoint. Sorting names into quartiles by at-midpoint share
(Table~\ref{tab:signing_quartile}), all four classifiers agree where few prints are at the
midpoint (Q1, $15\%$ share: $+0.74$ to $+1.80$); where most are (Q4, $91\%$), only the
abstaining rule retains a positive value while every force-signing rule goes to zero or
negative.

\begin{table}[htbp]
\centering
\begin{tabular}{lccccc}
\toprule
\textbf{At-mid quartile} & \textbf{At-mid share} & \textbf{Abstain} & \textbf{Tick} & \textbf{EMO} & \textbf{CLNV} \\
\midrule
Q1 (fewest at-mid) & $15\%$ & $+1.80$ & $+1.47$ & $+0.74$ & $+1.28$ \\
Q2 & $45\%$ & $+1.42$ & $+0.69$ & $+0.19$ & $+0.45$ \\
Q3 & $75\%$ & $+1.98$ & $+0.18$ & $-0.31$ & $+0.09$ \\
Q4 (most at-mid) & $91\%$ & $+2.74$ & $-0.27$ & $-0.52$ & $-0.31$ \\
\bottomrule
\end{tabular}
\caption{Three-minute footprint (bps) by classifier across quartiles of a name's
at-midpoint print share. Classifier disagreement appears only as the at-midpoint share
rises, locating it in the treatment of midpoint prints rather than in the underlying
measurement. The Q4 row carries a composition caveat given in the text.}
\label{tab:signing_quartile}
\end{table}

The Q4 cell requires a caveat that we state rather than bury, because it is the cell that
most favours our preferred convention. The abstaining rule's dependent variable is estimated
\emph{only} on non-midpoint prints, so the sorting variable mechanically governs the size of
the sample on which the outcome is measured. Across quartiles the number of signed prints
per name-day falls from $436$ in Q1 to $34$ in Q4---a factor of thirteen. The $+2.74$ in Q4
is therefore estimated from roughly $9\%$ of those names' hidden prints, and is
correspondingly noisier and potentially selected: abstention may be isolating the most
informative subsample rather than measuring a population. This does not affect the
within-name-day bifurcation of \S\ref{sec:decomp}, which does not sort on this variable, but
it does mean the quartile gradient should not be read as a clean cross-sectional
relationship.

\subsection{Two populations}
\label{sec:decomp}

Table~\ref{tab:decomp} measures the two populations separately, and is the paper's central
result. Hidden executions away from the midpoint are associated with a $+2.09$ bps
three-minute move that is flat across horizons ($+2.09 \to +2.24 \to +2.25$ over
$3/15/30$ minutes). Hidden executions at the midpoint carry a small \emph{negative} markout
($-0.42$ at three minutes, drifting to $-0.71$ at thirty), which is the signature of
uninformed liquidity provision absorbing transient price moves.

\begin{table}[htbp]
\centering
\begin{tabular}{lccc}
\toprule
\textbf{Hidden-execution subset} & \textbf{3 min} & \textbf{15 min} & \textbf{30 min} \\
\midrule
Aggressive (away from the midpoint) & $+2.09\;(t{=}25.3)$ & $+2.24\;(t{=}12.2)$ & $+2.25\;(t{=}13.3)$ \\
At the midpoint (tick-signed) & $-0.42\;(t{=}{-}6.6)$ & $-0.64\;(t{=}{-}7.0)$ & $-0.71\;(t{=}{-}10.7)$ \\
\bottomrule
\end{tabular}
\caption{Signed midpoint markout (bps) of aggressive versus at-midpoint hidden executions
(474 names, 2023--2024, $199{,}627$ name-days). Pooling the two is what produces the lower,
classifier-dependent aggregate of Table~\ref{tab:signing}.}
\label{tab:decomp}
\end{table}

We flag a definitional caution the reader should carry forward. Prints away from the
midpoint are, by construction, prints with an identifiable aggressor, so the statement
``aggressor-signed prints carry adverse selection and unsignable prints do not'' is partly a
property of trade classification rather than a purely economic finding. What is \emph{not}
definitional is the sign of the second population. A classification artifact would place
at-midpoint prints near zero; they instead carry a reliably negative markout that grows with
horizon.

That defence is only as good as the sign it rests on, and the at-midpoint row is tick-signed.
The tick rule labels a print a buy precisely when the price last ticked up, so if
short-horizon reversion follows that preceding tick, a negative forward markout would arise
mechanically for \emph{any} set of moments signed this way---no liquidity-provision
economics required, and the growth with horizon would be equally expected. Since this is the
load-bearing evidence that the bifurcation is economic rather than classificational, we test
it directly in \S\ref{sec:tickplacebo} rather than argue it.

The reading is that concealment is informationally bifurcated. Concealed \emph{aggressive}
intent is associated with adverse selection; concealed \emph{passive} midpoint provision is
not. This refines rather than restates the finding that hidden orders are informative: only
the aggressive half is, and reporting a single pooled number averages an informed
population against an uninformed one.

\subsection{The at-midpoint markout is not a tick-rule artifact}
\label{sec:tickplacebo}

The concern of the preceding paragraph has a sharp implication: if tick-rule conditioning
manufactures negative markouts, it should do so wherever it is applied, including at moments
that are not hidden midpoint executions at all. We test it two ways on the same panel.

\textbf{A matched placebo.} We take the at-midpoint print times of each name-day and shift
them all by a single random per-day offset. Inter-arrival structure---and therefore the burst
clustering, which is spacing-dependent---is preserved exactly, while the moments themselves
are decoupled from actual hidden midpoint executions. We then read the midpoint at those
shifted times, sign it by the identical tick rule, cluster identically, and measure the same
markouts. If the negative row is a property of tick-signing rather than of midpoint
concealment, the placebo should reproduce it.

\textbf{Tick-conditioning a known-informed population.} We re-sign the \emph{aggressive}
prints by the pure tick rule, discarding their quote-rule sign. This asks what tick
conditioning does to a population we have independent reason to believe is informed.

\begin{table}[htbp]
\centering
\begin{tabular}{lccc}
\toprule
\textbf{Tick-rule placebo} & \textbf{3 min} & \textbf{15 min} & \textbf{30 min} \\
\midrule
At-midpoint prints (tick-signed) & $-0.411\;(t{=}{-}9.0)$ & $-0.587\;(t{=}{-}10.4)$ & $-0.674\;(t{=}{-}11.1)$ \\
Matched placebo (shifted times, tick-signed) & $-0.022\;(t{=}{-}1.2)$ & $-0.088\;(t{=}{-}3.3)$ & $-0.110\;(t{=}{-}2.9)$ \\
\midrule
\multicolumn{4}{l}{\emph{Tick-conditioning the aggressive leg (3 min)}} \\
\quad Aggressive, quote-rule signed & \multicolumn{3}{l}{$+1.232\;(t{=}{+}26.5)$} \\
\quad Aggressive, pure tick-rule signed & \multicolumn{3}{l}{$-0.064\;(t{=}{-}1.6)$} \\
\bottomrule
\end{tabular}
\caption{Tick-rule placebo tests (470 names, 2023--2024, $235{,}940$ name-days). The placebo
recovers $5.4\%$ of the at-midpoint markout at three minutes and is statistically zero;
tick-signing the aggressive population destroys its footprint rather than reversing it.
Levels in the lower panel are estimated on this run's filter (which additionally requires
$180$s of within-session history and a valid pre-print quote) and are therefore not directly
comparable to Table~\ref{tab:decomp}'s $+2.09$; the informative quantity is the
\emph{within-panel} contrast between the two rows.}
\label{tab:tickplacebo}
\end{table}

Neither test supports the artifact reading. The placebo earns $-0.022$ bps at three minutes
against the at-midpoint row's $-0.411$---$5.4\%$ of it, and statistically indistinguishable
from zero. At fifteen and thirty minutes the placebo does become significant, but recovers
only $15\%$ and $16\%$ of the corresponding at-midpoint values, so tick-rule conditioning
contributes a detectable but small negative drift at longer horizons and essentially none at
the three-minute horizon the paper's tables use.

The second test is the more informative of the two. Re-signing the aggressive prints by the
pure tick rule does not turn their footprint negative; it drives it to zero
($-0.064$, $t=-1.6$, against $+1.232$ under the quote rule). What tick conditioning does is
\emph{destroy} directional information, not manufacture the appearance of contrarian
trading. Both results point the same way: the negative markout on at-midpoint hidden prints
is a property of those executions, not of the classifier applied to them.

We state the residual limitation. These tests establish that the tick rule does not
manufacture the sign; they do not establish that it recovers the \emph{correct} sign for any
individual midpoint print, and given the second row above---where tick-signing a known
population yields nothing---the tick rule is evidently a weak instrument. The at-midpoint
markout should therefore be read as evidence that this population differs sharply from the
aggressive one, which is what the bifurcation claims, rather than as a precise estimate of
uninformed liquidity provision. Signing midpoint prints against a consolidated quote, where
a print at the NASDAQ midpoint may sit away from the NBBO midpoint and acquire a genuine
sign, would settle it; see \S\ref{sec:limits}.

\subsection{The bifurcation is not an averaging artifact}

A pooled term structure that is flat could in principle blend continuation in some names
against reversion in others. Table~\ref{tab:regime_ts} splits the aggressive-hidden term
structure by price quintile. It does not blend: the markout is flat or gently building in
every quintile and reverts in none. It is larger in cheap, tick-constrained names ($+3.96$
bps in Q1 against $+1.42$ in Q5), which Section~\ref{sec:adverse} shows is the spread
effect.

\begin{table}[htbp]
\centering
\begin{tabular}{lcccc}
\toprule
\textbf{Price quintile} & \textbf{Avg.\ price} & \textbf{3 min} & \textbf{15 min} & \textbf{30 min} \\
\midrule
Q1 (tick-constrained) & \$14 & $+3.96$ & $+4.21$ & $+4.19$ \\
Q2 & \$34 & $+1.89$ & $+2.62$ & $+2.61$ \\
Q3 & \$58 & $+1.53$ & $+1.51$ & $+1.46$ \\
Q4 & \$97 & $+1.53$ & $+1.58$ & $+1.43$ \\
Q5 (large-price) & \$281 & $+1.42$ & $+1.43$ & $+1.46$ \\
\bottomrule
\end{tabular}
\caption{Aggressive-hidden markout (bps) by price quintile (474 names, 2023--2024,
$\sim\!44{,}000$ name-days per quintile; all $t>12$). Flat or building in every quintile,
reverting in none.}
\label{tab:regime_ts}
\end{table}

\subsection{The intraday term structure, on a subsample}
\label{sec:termstructure}

The horizons above stop at thirty minutes. To instrument the rest of the session we extend
the aggressive-hidden markout to one hour, two hours and the burst-to-close midpoint, and
net each horizon against a \emph{time-of-day-stratified} placebo that draws each burst's
counterfactual time within its own 30-minute intraday bucket, so the null preserves the
U-shaped intensity and drift profile a uniform placebo would ignore.

\begin{table}[htbp]
\centering
\begin{tabular}{lcccc}
\toprule
\textbf{Aggressive-hidden markout} & \textbf{Burst (bps)} & \textbf{TOD placebo (bps)} & \textbf{Net (bps)} & \textbf{Net $t$ (day)} \\
\midrule
3 min    & $+2.023$ & $+0.007$ & $+2.017$ & $+28.3$ \\
15 min   & $+2.122$ & $-0.022$ & $+2.145$ & $+14.8$ \\
30 min   & $+2.166$ & $-0.101$ & $+2.268$ & $+15.5$ \\
1 hour   & $+2.131$ & $-0.145$ & $+2.283$ & $+15.3$ \\
2 hour   & $+2.076$ & $-0.289$ & $+2.378$ & $+10.3$ \\
to close & $+2.014$ & $-0.415$ & $+2.433$ & $+10.0$ \\
\bottomrule
\end{tabular}
\caption{Intraday term structure of the aggressive-hidden footprint on the \emph{full panel}
(470 names, 2023--2024, $\approx\!206{,}070$ name-days after the outlier rule). Each horizon
is netted against a time-of-day-stratified placebo drawing each burst's counterfactual time
within its own 30-minute bucket; $t$-statistics are Newey--West on the daily-mean net series.
The \emph{gross} markout is flat across the whole session---$+2.02$ at three minutes,
$+2.01$ measured to the close---so the profile does not depend on the placebo to stay level.}
\label{tab:hidden_term}
\end{table}

The profile is flat across the entire session on both measures. The gross markout runs
$+2.02$, $+2.12$, $+2.17$, $+2.13$, $+2.08$, $+2.01$ from three minutes to the close, and the
placebo is small at every horizon ($|\cdot| \le 0.42$), so the net rises only gently from
$+2.02$ to $+2.43$. A hidden burst's directional association with price is established within
roughly three minutes and then neither grows nor decays for the remaining six and a half
hours.

One structural caveat attaches to the final row and we state it rather than bury it. Every
other horizon measures a fixed elapsed time, whereas ``to close'' measures to a fixed
endpoint, so a counterfactual timestamp drawn earlier in its bucket has a systematically
longer window than the burst it replaces; this is why the placebo drifts monotonically more
negative with horizon. In the earlier 48-name version of this table that mattered, because
the gross to-close markout there fell to $+1.42$ and the row netted to $+2.08$ only because
the placebo was simultaneously more negative. On the full panel it does not: the gross
markout at the close is $+2.014$, statistically indistinguishable from its three-minute
value, so the flat profile holds without any placebo adjustment at all.

\subsection{The sign is not a stale-midpoint artifact}
\label{sec:signrobust}

Because trades are signed against a reconstructed midpoint, a print above a \emph{lagging}
midpoint could be mechanically classified a buy just before the midpoint ticks up,
manufacturing a positive markout. Table~\ref{tab:signrobust} re-signs three ways.

\begin{table}[htbp]
\centering
\begin{tabular}{lcc}
\toprule
\textbf{Signing rule (3-min aggressive-hidden markout)} & \textbf{bps} & \textbf{NW $t$} \\
\midrule
Baseline (midpoint at the trade) & $+2.00$ & $+23.8$ \\
Staler midpoint (quote 1s earlier) & $-0.02$ & $-0.4$ \\
Forward midpoint (quote 1s later; look-ahead) & $-1.59$ & $-39.2$ \\
Outside-the-quote prints only (unambiguous sign) & $+2.09$ & $+12.5$ \\
\bottomrule
\end{tabular}
\caption{Signing robustness (474 names, 2023--2024, $146{,}393$ name-days). Signing against
a staler midpoint degrades rather than inflates the footprint, the opposite of the
stale-midpoint prediction; and on prints executing outside the quote, where the sign cannot
be wrong, the footprint is undiminished.}
\label{tab:signrobust}
\end{table}

Signing against a deliberately staler midpoint does not inflate the footprint---it collapses
it to zero, because a staler midpoint misclassifies more prints, which is the opposite of
what the stale-midpoint hypothesis predicts. Most directly, restricting to prints executing
\emph{outside} the prevailing quote, where the aggressor's side is unambiguous regardless of
any midpoint error, leaves the footprint at $+2.09$ bps ($t=12.5$). The measurement is
strongest exactly where the sign is most certain.

We record here, and return to in Section~\ref{sec:overstate}, that this outside-the-quote
subset is also the one on which the pre-print evidence is cleanest---and there it cuts the
other way.

\section{The Footprint Is an Adverse-Selection Quantity}
\label{sec:adverse}

\subsection{It scales with the spread, and not with volatility}

If the aggressive footprint is the adverse-selection component of the spread, it carries a
sharp cross-sectional prediction: it should be a roughly constant \emph{fraction} of the
quoted half-spread across names, and should track the spread specifically rather than
volatility, illiquidity or price level generally. This is a stronger test than anything
that precedes it, because a mechanical or volatility-scaling account does not predict it.

The prediction holds. Across 469 names with at least 100 name-days, the correlation between
the quoted half-spread and the permanent price-impact component is $+0.85$, and the ratio of
the two is centered on $0.34$ with an interquartile range of $0.28$--$0.42$. Sorting into
quintiles by quoted half-spread, the spread ranges over a factor of five ($2.47$ to $12.54$
bps) while the adverse-selection \emph{share} moves only between $0.29$ and $0.38$.

\begin{table}[htbp]
\centering
\begin{tabular}{lccc}
\toprule
\textbf{Regressor (standardized)} & \textbf{(1)} & \textbf{(2)} & \textbf{(3)} \\
\midrule
Quoted half-spread & $+0.850$ & & $+0.930$ \\
                   & $(+34.8)$ & & $(+25.2)$ \\
Realized volatility & & $+0.573$ & $+0.008$ \\
                    & & $(+15.0)$ & $(+0.2)$ \\
$\log$ price & & & $-0.042$ \\
             & & & $(-0.9)$ \\
$\log$ dollar volume & & & $+0.224$ \\
                     & & & $(+4.8)$ \\
\midrule
$R^2$ & $0.723$ & $0.328$ & $0.755$ \\
\bottomrule
\end{tabular}
\caption{Cross-name determinants of the permanent price-impact component (465 names,
2023--2024; all variables standardized, $t$-statistics in parentheses). Univariately both
the quoted half-spread and realized volatility predict the footprint. Jointly, the spread
coefficient is undiminished while volatility falls to zero.}
\label{tab:scaling}
\end{table}

Column (3) is the discriminating test. Realized volatility predicts the footprint strongly
on its own ($+0.573$, $t=15.0$). Entered jointly with the quoted half-spread it collapses to
$+0.008$ ($t=0.2$) while the spread coefficient is if anything larger. The footprint is not
a volatility phenomenon that happens to appear in wide-spread names; it is a spread
phenomenon for which volatility was proxying.

One caveat constrains this. Wide-spread names have a coarser effective price grid, so a
given midpoint move is mechanically larger there, and that channel inflates measured impact
in wide-spread names in a way we cannot fully separate from the informational one. It does
not, however, explain the \emph{stability of the ratio} across quintiles, nor the
disappearance of the volatility coefficient under the spread control, since tick coarseness
and volatility are themselves positively related.

\subsection{Its share of trading costs}

Table~\ref{tab:spread} places the footprint on the \citet{huangstoll1997} accounting,
estimated on the aggressive-hidden trades with $\tau = 180$ seconds.

\begin{table}[htbp]
\centering
\begin{tabular}{lcc}
\toprule
\textbf{Huang--Stoll component (aggressive-hidden trades)} & \textbf{bps} & \textbf{NW $t$} \\
\midrule
Quoted half-spread & $+6.39$ & $+27.4$ \\
Effective half-spread (taker pays) & $+4.00$ & $+23.5$ \\
\quad = Price impact (adverse selection) & $+2.02$ & $+32.4$ \\
\quad + Realized half-spread (provider retains) & $+1.97$ & $+16.2$ \\
\bottomrule
\end{tabular}
\caption{Effective/realized spread decomposition \citep{huangstoll1997} for
aggressive-hidden trades (474 names, 2023--2024, $203{,}142$ name-days). Reported for the
\emph{ratio} it yields, not as independent confirmation---see the text.}
\label{tab:spread}
\end{table}

We are explicit about what this does and does not add, because it is easy to overstate. The
Huang--Stoll price-impact component is defined as $D_t(M_{t+\tau} - M_t)$, which is
\emph{algebraically the markout itself}; the realized half-spread is then the effective
half-spread minus that quantity. The near-coincidence of the $+2.02$ bps component with the
$+2.09$ bps markout of Table~\ref{tab:decomp} is therefore not independent corroboration.
It is the same estimand computed twice on slightly different name-day filters, and any
claim that the decomposition \emph{confirms} the markout would be circular.

What the decomposition genuinely contributes is the ratio, which the markout alone cannot
supply: the permanent component is $2.02$ of a $4.00$ bps effective half-spread, or $51\%$
of what the taker actually pays, and $32\%$ of the $6.39$ bps quoted half-spread. Roughly
half the effective cost of aggressive hidden execution is compensation for information and
roughly half is the provider's gross margin.

\subsection{It is incremental to visible order-flow imbalance}
\label{sec:ofi}

Hidden flow might simply proxy for the visible order-flow imbalance already known to
forecast short-horizon returns \citep{cont2014}, in which case the footprint would be a
re-description rather than an addition. We bin each name-day on a fixed ten-second clock
and regress the forward three-minute midpoint return on quote-rule-signed aggressive-hidden
volume and on the visible Cont--Kukanov--Stoikov imbalance, both standardized within
name-day, univariately and jointly.

\begin{table}[htbp]
\centering
\begin{tabular}{lcc}
\toprule
\textbf{Predictor of forward 3-min midpoint return} & \textbf{bps per SD} & \textbf{NW $t$} \\
\midrule
\multicolumn{3}{l}{\emph{Univariate}} \\
\quad Aggressive-hidden signed volume & $+0.160$ & $+12.2$ \\
\quad Visible OFI \citep{cont2014} & $+2.351$ & $+54.6$ \\
\midrule
\multicolumn{3}{l}{\emph{Joint}} \\
\quad Aggressive-hidden signed volume & $+0.198$ & $+22.9$ \\
\quad Visible OFI & $+2.330$ & $+54.2$ \\
\bottomrule
\end{tabular}
\caption{Hidden flow versus visible order-flow imbalance (474 names, 2023--2024,
$168{,}556$ name-days, ten-second bins; day-level Fama--MacBeth with Newey--West errors).
Controlling for visible OFI does not attenuate the hidden-flow coefficient.}
\label{tab:ofi}
\end{table}

The hidden-flow coefficient is not attenuated by the control; it \emph{rises}, from $+0.160$
to $+0.198$ bps per SD, with the $t$-statistic nearly doubling. This requires explanation,
because two positively correlated regressors that both positively predict the outcome should
attenuate one another. Enhancement instead indicates classical suppression, which arises
when the conditional correlation between regressors is negative---and here it is. The
mechanism is specific to this pair: a hidden \emph{buy} consumes displayed depth at the ask,
and ask-side depletion enters the Cont--Kukanov--Stoikov imbalance with a \emph{negative}
sign, so the two carry mildly opposing contemporaneous components even though both forecast
returns positively.

The magnitude reconciles exactly, which is worth showing because a correlation this small
looks too weak to move a coefficient by a quarter. For standardized regressors,
$b_H^{\text{joint}} = (b_H^{\text{uni}} - \rho\, b_O^{\text{uni}})/(1-\rho^2)$, and
substituting the measured values gives
$(0.1601 - (-0.0161)(2.3508))/(1-0.0161^2) = 0.1980$ against the $0.1981$ reported. The
suppression term is $24\%$ of $b_H^{\text{uni}}$ precisely because $b_O^{\text{uni}}$ is
$14.7$ times larger: a small correlation multiplied by a much larger coefficient is a large
adjustment \emph{relative to the smaller one}. The implied $\rho$ of $-0.016$ also matches
the directly measured median name-day correlation of $-0.027$ on the same clock.

Visible OFI is by far the larger short-horizon predictor in absolute terms, an order of
magnitude above the hidden coefficient, and we do not claim otherwise. The point is only
that the hidden footprint survives it intact.

\section{Two Non-Informational Accounts}
\label{sec:alternatives}

\subsection{Information events}

The most natural alternative to an adverse-selection reading is that aggressive hidden
executions cluster around information events---earnings, guidance, macro releases---and that
price moves because of the announcement, with the execution a correlate. We obtain quarterly
earnings dates from Yahoo Finance, retrieved through the \texttt{yfinance} interface, for the
418 panel names that report ($4{,}925$ announcements over 2023--2024; the remaining 56 are
ETFs, which do not). We proxy for \emph{unscheduled} news by the magnitude of the name's
daily return. We note that Yahoo Finance is a retail-grade source for announcement
timestamps; a commercial calendar would be preferable, and the dates enter only as an
exclusion window rather than as a treatment variable.

\begin{table}[htbp]
\centering
\begin{tabular}{lccc}
\toprule
\textbf{Sample} & \textbf{Footprint (bps)} & \textbf{NW $t$} & \textbf{name-days} \\
\midrule
All name-days (earnings-covered names) & $+1.910$ & $+29.3$ & $154{,}577$ \\
\midrule
\quad Earnings window ($\pm1$ day) & $+3.002$ & $+13.9$ & $7{,}657$ \\
\quad Excluding earnings window & $+1.874$ & $+29.6$ & $146{,}920$ \\
\quad Excluding earnings ($\pm3$ days) & $+1.858$ & $+31.4$ & $141{,}715$ \\
\midrule
\quad Largest $10\%$ of daily moves & $+2.959$ & $+27.1$ & $17{,}559$ \\
\quad Excluding largest $10\%$ of moves & $+1.751$ & $+23.6$ & $143{,}527$ \\
\midrule
\textbf{Excluding earnings $\pm2$d \emph{and} top $5\%$ moves} & $\mathbf{+1.760}$ & $\mathbf{+26.8}$ & $143{,}111$ \\
\bottomrule
\end{tabular}
\caption{The footprint after removing information-event windows (474-name panel,
2023--2024; three-minute quote-rule markout). On the joint-clean subsample---no earnings
window, no top-decile move, $88\%$ of all name-days---the footprint retains $94\%$ of its
value.}
\label{tab:news}
\end{table}

The confound is not supported. Excluding earnings windows moves the footprint from $+1.910$
to $+1.874$; excluding the largest decile of daily moves leaves $+1.751$; excluding both,
retaining $88\%$ of the panel, leaves $+1.760$ against $+1.878$, or $94\%$ retention with
the $t$-statistic essentially unchanged. The same pattern holds on the Huang--Stoll series
($+2.02 \to +1.99$).

The event cells are themselves informative rather than a nuisance: the footprint roughly
doubles around earnings ($+3.00$) and rises monotonically with the size of the daily move
($+2.96$, $+3.43$, $+3.98$ at the top decile, ventile and percentile). That is what adverse
selection predicts when more private information is in play, and it is difficult to
reconcile with a purely mechanical account of the same quantity.

We note the limit. Return magnitude is a coarse proxy for news: it captures unscheduled
information only insofar as it moved prices that day, and cannot distinguish news arrival
from a liquidity shock. A timestamped intraday news feed would be the full treatment. What
the test establishes is that the scheduled-announcement channel accounts for essentially
none of the measured footprint.

\subsection{Mechanical quote displacement}
\label{sec:depletion}

The more serious competing account is mechanical. An aggressive execution consumes displayed
depth; if the level is depleted the best quote moves, the midpoint moves with it, and if the
level is never replenished the displacement persists. Such an account reproduces every
stylized fact above---a footprint appearing instantly, never reverting, scaling with the
spread, largest in coarse-tick names---with no information conveyed. Ratio stability across
spread quintiles does not discriminate, being equally a prediction of queue geometry.

The test must therefore condition on something known strictly \emph{before} the print. We
use the pre-trade depletion ratio $d = (\text{trade size}) / (\text{displayed depth at the
touch, immediately pre-trade})$, both components observed before any price response. The
mechanical account makes a sharp prediction about $d$: displacement requires consuming the
queue, so a print taking a small slice of a deep book cannot move the midpoint. The
informational account makes no such prediction.

\begin{table}[htbp]
\centering
\begin{tabular}{lccr}
\toprule
\textbf{Pre-trade depletion ratio $d$} & \textbf{Markout (bps)} & \textbf{NW $t$} & \textbf{prints} \\
\midrule
Q1 (smallest slice of the queue) & $+1.169$ & $+25.0$ & $38.4$M \\
Q2 & $+1.168$ & $+26.3$ & $38.2$M \\
Q3 & $+1.289$ & $+28.3$ & $42.8$M \\
Q4 & $+1.381$ & $+28.5$ & $33.9$M \\
Q5 (largest) & $+1.426$ & $+35.7$ & $36.9$M \\
\midrule
\textbf{$d < 0.10$: consumed $<\!10\%$ of displayed depth} & $\mathbf{+1.134}$ & $\mathbf{+26.3}$ & $46.0$M \\
$d \geq 1.0$: consumed the entire visible level & $+1.788$ & $+37.2$ & $62.0$M \\
\midrule
All aggressive prints & $+1.303$ & $+34.3$ & \\
\bottomrule
\end{tabular}
\caption{Three-minute markout by pre-trade depletion ratio (470 names, 2023--2024,
$127{,}667$ name-days). Both trade size and displayed depth are measured before the print,
so the conditioning variable cannot be contaminated by the price response.}
\label{tab:depletion}
\end{table}

Prints consuming less than $10\%$ of displayed depth---$24\%$ of aggressive prints, 46
million of them---carry $+1.134$ bps at $t=26.3$, which is $87\%$ of the unconditional
average. These trades leave the touch essentially intact, so there is no displacement to
explain the subsequent move, yet the move is nearly as large as for prints that sweep the
level entirely. The quintile gradient tells the same story: $d$ rises by orders of magnitude
from Q1 to Q5 while the footprint rises only $22\%$, whereas a mechanical account predicts
impact roughly proportional to the fraction of queue consumed, which would put Q1 near zero.

We do not read this as excluding a mechanical component. The gradient is monotone and the
$d \geq 1.0$ cell is $58\%$ larger than the $d<0.10$ cell, consistent with genuine
displacement adding to an informational base, and with larger orders being more informed.

A scope note on the construction is essential here, and it limits the test more than we first
allowed. A type-5 print executes against non-displayed resting liquidity, so when the hidden
order rests inside the spread the print does not itself consume the displayed queue; $d$ is a
normalization of trade size by available visible depth, not a literal fraction of the queue
consumed. A print with $d<0.10$ therefore has \emph{not} been shown to displace nothing---it
has been shown to be small relative to visible depth, which is compatible with the same
aggressor's marketable remainder sweeping the touch microseconds later. The mechanical
channel that actually applies to hidden prints is that sequential sweep, and
Table~\ref{tab:depletion} does not test it.

\subsection{The sequential-sweep channel}
\label{sec:sweep}

We therefore observe the sweep directly. For each aggressive hidden print we record the touch
on the side being taken immediately before the print and again $100$ms later, and call the
level \emph{swept} if its price moved away or its size more than halved at an unchanged price.
The channel is common: $29.8\%$ of aggressive hidden prints are followed by a sweep within
$100$ms.

The measurement requires care, because the obvious test is circular. Conditioning on the
touch moving inside the first second and then measuring a markout that begins at the print
conditions on most of the outcome---the footprint is largely impounded within that first
second, and for a buy a rising ask \emph{is} a rising midpoint. We therefore separate the
conditioning window from the measurement window: prints are classified on $[t, t+100\text{ms}]$
alone, and the forward move is measured from $t+1$s and from $t+5$s, strictly after the
displacement. Table~\ref{tab:sweep} reports both the circular and the clean versions so the
contamination is visible rather than asserted.

\begin{table}[htbp]
\centering
\begin{tabular}{lccc}
\toprule
\textbf{Markout window} & \textbf{Swept} & \textbf{Unswept} & \textbf{Gap} \\
\midrule
From $t$ to $t+180$s \emph{(circular; reference only)} & $+3.271\;(58.6)$ & $+0.464\;(22.0)$ & $+2.807$ \\
\midrule
From $t+1$s to $t+180$s & $+0.374\;(9.7)$ & $+0.439\;(22.0)$ & $-0.066$ \\
From $t+5$s to $t+185$s & $+0.203\;(5.8)$ & $+0.260\;(15.1)$ & $-0.057$ \\
\bottomrule
\end{tabular}
\caption{Forward move by whether the touch was swept within $100$ms of the print (470 names,
2023--2024, $235{,}940$ name-days; bps, NW $t$ in parentheses; median $153$ swept and $378$
unswept prints per name-day). The first row conditions on a quote move inside the window it
then measures; the lower two do not. The $18\times$ split in the first row is absent from
both clean rows.}
\label{tab:sweep}
\end{table}

Two findings, and the second is uncomfortable for the reading advanced above.

\textbf{The sweep does not mark differentially informative flow.} Measured from after the
displacement, swept and unswept prints are indistinguishable, and if anything swept prints
carry slightly \emph{less} forward move ($+0.374$ against $+0.439$ from $t+1$s; $+0.203$
against $+0.260$ from $t+5$s). Whatever a sweep signals, it does not predict a larger
subsequent revision. The apparent $18\times$ concentration is entirely an artifact of the
overlap between conditioning and measurement.

\textbf{But the level of the measured footprint is concentrated in swept prints, and for
those prints it coincides with the displacement.} Weighting the from-$t$ row by frequency,
$29.8\% \times 3.271 + 70.2\% \times 0.464 = 1.30$ bps, reproducing the all-print average of
Table~\ref{tab:depletion}; of that, $75\%$ comes from the swept minority. And for swept
prints almost the whole of it occurs inside the first second ($+3.271$ from $t$ against
$+0.374$ from $t+1$s), which is exactly when the touch moves. We therefore withdraw the
claim, made in an earlier version of this section, that displacement cannot account for the
level of the footprint. On the direct test it accounts for a majority of it.

What this does \emph{not} establish is that the footprint is uninformative, and the
distinction matters. An aggressor whose order exceeds displayed depth sweeps the touch, and
an aggressor who wants more liquidity than is displayed is plausibly one who knows something;
under the Glosten--Milgrom mechanism the quote moving as informed flow takes liquidity is the
price-discovery process, not an artifact of it. ``Mechanical displacement'' and ``information
impounded through the trade'' are the same event at a $100$-millisecond resolution, and
nothing in this design separates them. What survives cleanly is the unswept population:
$70\%$ of aggressive hidden prints, displacing nothing at the touch, carrying $+0.464$ bps
from the print and $+0.439$ from a second after it. That is a real but small footprint, and
it is the most conservative reading of the level our data supports.

\section{What Unconditional Markouts Overstate}
\label{sec:overstate}

Everything above is a forward markout: a measurement of where price goes from the print. A
forward markout cannot distinguish response from continuation. If conditioning on aggressive
same-side execution selects moments at which the midpoint is already travelling in that
direction, the markout measures a move already underway. This section tests that directly
and reports that it is substantially true.

\subsection{The event-time profile}

Table~\ref{tab:preprint} measures the signed midpoint move over windows \emph{ending} at the
print as well as the forward markout.

\begin{table}[htbp]
\centering
\begin{tabular}{lcc}
\toprule
\textbf{Signed midpoint move (bps)} & \textbf{All aggressive} & \textbf{Outside the quote} \\
\midrule
$-180$s $\rightarrow$ $t$ & $-1.541$ $(-60.6)$ & --- \\
$-30$s $\rightarrow$ $t$ & $-1.715$ $(-56.7)$ & $+1.195$ $(+35.6)$ \\
$-5$s $\rightarrow$ $t$ & $-1.783$ $(-50.2)$ & $+0.965$ $(+23.6)$ \\
\midrule
$t$ $\rightarrow$ $+180$s (the markout) & $+1.159$ $(+35.2)$ & $+1.322$ $(+27.3)$ \\
\bottomrule
\end{tabular}
\caption{Event-time profile around aggressive hidden prints (468 names, $206{,}235$
name-days; NW $t$-statistics in parentheses). Moves are signed by trade direction, so a
positive pre-print figure means price was already travelling \emph{with} the trade before it
arrived. The two columns disagree in sign; the outside-the-quote column is the informative
one.}
\label{tab:preprint}
\end{table}

\textbf{The pooled column is contaminated and should not be read.} On all aggressive prints
the pre-print drift is strongly negative, appearing to show price moving against the trade
before it arrives. This is an artifact of sign assignment. A print is classified a buy when
it executes above the \emph{contemporaneous} midpoint, so a print at an unchanged price is
mechanically labelled a buy whenever the midpoint has just fallen. Falling-midpoint-then-buy
is built into the classification, and the negative pre-drift is the stale-midpoint problem
of Section~\ref{sec:signrobust} reappearing in event time rather than evidence of contrarian
trading.

\textbf{The outside-the-quote column reverses it.} On prints executing beyond the prevailing
quote, where the aggressor's identity cannot be mistaken, price moves \emph{with} the trade
direction before it arrives: $+1.195$ bps over the preceding thirty seconds ($t=35.6$) and
$+0.965$ over the preceding five. Over the window $[-30\text{s},+180\text{s}]$ the signed
move totals $+2.52$ bps, of which $47.5\%$ occurs before the print.

\subsection{Is the markout continuation of the pre-drift?}
\label{sec:decomposition}

The accounting above shows that roughly half the signed move precedes the print. It does
\emph{not} show that the forward markout is continuation of that pre-drift, and the two are
different claims: the first is about where the move sits in event time, the second about
whether the pre-drift predicts the post-print move. We test the second directly. Within each
name-day we regress the signed forward markout on the signed pre-print drift over the
preceding thirty seconds,
\begin{equation}
m_h \;=\; a_h \;+\; b_h\, d \;+\; \varepsilon,
\qquad\text{so that}\qquad
\overline{m_h} \;=\; a_h \;+\; b_h\, \overline{d}
\label{eq:decomp}
\end{equation}
holds exactly within the name-day. The intercept $a_h$ is the component of the markout
orthogonal to the prior price path---the quantity a path-conditioning estimator such as the
VAR is designed to isolate---and $b_h\overline{d}$ is the part attributable to continuation.
Averaging both across the panel decomposes the reported markout into two pieces that sum to
it by construction.

\begin{table}[htbp]
\centering
\begin{tabular}{lccc}
\toprule
& \textbf{3 min} & \textbf{10 min} & \textbf{30 min} \\
\midrule
\multicolumn{4}{l}{\emph{Outside-the-quote prints} (mean pre-drift $\overline{d} = +1.301$, $t=+37.0$)} \\
\quad Markout $\overline{m_h}$ & $+1.379\;(32.0)$ & $+1.228\;(23.4)$ & $+1.067\;(16.3)$ \\
\quad Orthogonal to pre-drift, $a_h$ & $+1.407\;(36.0)$ & $+1.285\;(24.2)$ & $+1.154\;(17.3)$ \\
\quad Continuation, $b_h\overline{d}$ & $-0.028\;(-1.9)$ & $-0.057\;(-3.1)$ & $-0.088\;(-3.2)$ \\
\midrule
\multicolumn{4}{l}{\emph{All aggressive prints} (mean pre-drift $\overline{d} = -1.684$, $t=-55.6$)} \\
\quad Markout $\overline{m_h}$ & $+1.153\;(36.5)$ & $+1.099\;(24.8)$ & $+1.030\;(16.9)$ \\
\quad Orthogonal to pre-drift, $a_h$ & $+1.005\;(28.9)$ & $+0.885\;(21.5)$ & $+0.756\;(11.9)$ \\
\quad Continuation, $b_h\overline{d}$ & $+0.148\;(7.8)$ & $+0.215\;(9.7)$ & $+0.274\;(13.8)$ \\
\bottomrule
\end{tabular}
\caption{Decomposition~\eqref{eq:decomp} of the forward markout into a component orthogonal
to the preceding thirty seconds of price movement and a continuation component (470 names,
2023--2024; $119{,}483$ name-days outside-the-quote, $201{,}271$ all-aggressive; bps, NW $t$
in parentheses). Rows sum to the markout by construction. On the clean outside-the-quote
population the continuation component is economically zero and slightly \emph{negative}: the
pre-drift does not forecast the post-print move.}
\label{tab:decomposition}
\end{table}

The result is not what the event-time accounting led us to expect. On the outside-the-quote
population---the one where the sign cannot be wrong---price does drift $+1.301$ bps into the
print, but that drift carries essentially none of the forward markout. The continuation
component is $-0.028$ bps at three minutes, $-2\%$ of the total and of the wrong sign, and
reaches only $-0.088$ at thirty. The slope $b_h$ is mildly negative throughout: prints
arriving after a \emph{larger} run-up are followed by slightly \emph{smaller} forward moves,
not larger ones. Essentially the entire markout is orthogonal to the prior path. The
all-aggressive column tells the same story with the sign of $\overline{d}$ flipped by the
contamination described above, and a continuation share of $13$--$27\%$.

Two conclusions follow, and they pull in opposite directions.

\textbf{The specific reverse-causality mechanism is not supported.} ``The markout measures
continuation of a move already underway'' makes a testable prediction---that conditioning on
the move already underway should absorb the markout---and the prediction fails. Roughly half
the signed move does precede the print, but the two halves are not the same phenomenon: the
pre-drift is uninformative about what follows.

\textbf{This does not restore a causal reading, and it does not reconcile the VAR.} Ruling
out one confound is not identification: flow may still be selected on information the
observable price path does not reveal, and \eqref{eq:decomp} conditions on a single
thirty-second summary of that path rather than on its full history. We therefore retain the
association language of \S\ref{sec:intro}.

\subsection{Consequences}

First, we make no causal claim. What we report is an association between aggressive hidden
execution and contemporaneous price discovery, with the direction of causality unresolved.

Second---and we correct an earlier reading of our own evidence here---the pre-print drift
does \emph{not} reconcile the markout with the structural benchmark. Table~\ref{tab:hasbrouck}
reports a \citet{hasbrouck1991} permanent--transitory VAR, re-estimated on one-minute bins
with 30 lags so that the horizons of interest lie inside the model's memory rather than
beyond it, with ridge shrinkage in place of discarding explosive fits and all name-days
retained rather than conditioning on stationarity. The cumulative response peaks within a
minute, is barely distinguishable from zero at three, and turns significantly negative by
ten and thirty. The natural reading---that the VAR conditions on the prior return path while
the markout does not, so the gap is the pre-drift---is the one Table~\ref{tab:decomposition}
rules out: conditioning on the prior path removes essentially none of the markout. What
separates the two estimators is therefore specification, not contamination, and
Table~\ref{tab:vargrid} identifies which feature of the specification is responsible.

\begin{table}[htbp]
\centering
\begin{tabular}{lccc}
\toprule
\textbf{Cumulative response to a 1-SD flow innovation} & \textbf{bps} & \textbf{NW $t$} & \textbf{95\% CI} \\
\midrule
1 minute  & $+0.031$ & $+7.80$ & $[+0.024, +0.039]$ \\
3 minutes & $+0.015$ & $+2.27$ & $[+0.002, +0.029]$ \\
10 minutes & $-0.046$ & $-4.17$ & $[-0.068, -0.025]$ \\
30 minutes & $-0.195$ & $-6.57$ & $[-0.253, -0.137]$ \\
\bottomrule
\end{tabular}
\caption{Re-estimated \citet{hasbrouck1991} VAR: one-minute bins, 30 lags (30 minutes of
model memory), ridge shrinkage, all name-days retained (473 names, $121{,}252$ name-days).
The stationary-only subsample gives the same pattern ($-0.049$ at 10 minutes, $-0.160$ at
30). Magnitudes are not comparable to the markout in levels---ridge shrinkage attenuates the
coefficients and the response is denominated per flow-innovation standard deviation on a
one-minute clock---so we read the sign and shape, not the scale.}
\label{tab:hasbrouck}
\end{table}

Two features distinguish the current specification from the earlier one that reported a
permanent response: the sampling clock together with the induced memory horizon, and ridge
shrinkage. Table~\ref{tab:vargrid} crosses them.

\begin{table}[htbp]
\centering
\begin{tabular}{lcccc}
\toprule
\textbf{Specification} & \textbf{3 min} & \textbf{10 min} & \textbf{30 min} & \textbf{\% stat.} \\
\midrule
10s clock, 12 lags, no ridge & $+0.043\;(10.0)$ & $+0.047\;(9.4)$ & $+0.047\;(9.0)$ & $99\%$ \\
10s clock, 12 lags, ridge & $+0.042\;(9.9)$ & $+0.046\;(9.6)$ & $+0.048\;(10.6)$ & $99\%$ \\
60s clock, 30 lags, no ridge & $+0.017\;(2.6)$ & $-0.047\;(-4.1)$ & $-0.145\;(-9.5)$ & $88\%$ \\
60s clock, 30 lags, ridge & $+0.015\;(2.4)$ & $-0.051\;(-4.8)$ & $-0.149\;(-9.4)$ & $89\%$ \\
\bottomrule
\end{tabular}
\caption{Cumulative VAR response to a one-SD flow innovation (bps) across the $2\times2$
specification grid (470 names, 2023--2024; NW $t$ in parentheses; responses trimmed at
$|\cdot|\le50$ bps, since non-stationary fits on the one-minute clock have heavy tails---at a
$10^4$ trim the three-minute and ten-minute columns are unchanged while the thirty-minute
column becomes noise, $t$ between $-1.9$ and $+1.4$). Rows 1--2 and 3--4 are within a
rounding error of one another; rows 1 and 3 differ in sign.}
\label{tab:vargrid}
\end{table}

The comparison is unambiguous. Ridge shrinkage does nothing: holding the clock fixed, the
responses agree to the third decimal at every horizon. The clock does everything: at ten and
thirty minutes the one-minute specification is negative and the ten-second specification is
positive.

The reason is visible in the first row. A VAR(12) on a ten-second clock carries $120$ seconds
of memory, so responses requested at three, ten and thirty minutes all lie beyond the model's
horizon and return its asymptote---which is why that row reports $+0.043$, $+0.047$, $+0.047$,
the same number three times. The apparent permanence of the original specification was not an
estimate of behaviour at ten and thirty minutes; it was the model's fixed point, reported at
horizons it could not see. Only the one-minute clock with thirty lags places those horizons
inside the estimable range, and there the response decays and turns negative.

This resolves what the specification difference is, and it removes the earlier estimate as
evidence for permanence. It does not resolve the disagreement with the markouts, which
remains: a path-conditioning estimator with genuine memory finds decay where the direct
measurement finds none, and \S\ref{sec:decomposition} has already shown that the obvious
bridge between them---continuation of pre-print drift---does not carry the weight. We report
the VAR as a negative result, do not cite it as corroborating permanence, and leave the
disagreement open.

Third, the post-print component is not zero and is not inflated by the pre-drift: $+1.379$
bps at three minutes on the clean population, of which $+1.407$ is orthogonal to the
preceding thirty seconds. Something happens after the print that does not happen before it,
and it is not a continuation of what happened before it.

Establishing that it is \emph{caused} by the print would require variation in flow unrelated
to the information environment, which this design does not provide;
Table~\ref{tab:decomposition} removes one specific alternative rather than delivering
identification. We record the causal question as open, and note that it does \emph{not} bear
on the bifurcation of Section~\ref{sec:bifurcation}, which compares two contemporaneous
populations rather than claiming what an execution does to prices.

\section{Scope and Limits}
\label{sec:limits}

\subsection{One venue}

This is the binding limitation and we state it first. LOBSTER reconstructs the NASDAQ book
alone, so every midpoint, quote and depth here is NASDAQ's rather than the consolidated
NBBO. The pre-print evidence of Section~\ref{sec:overstate} makes this first-order rather
than incidental: $+1.195$ bps of pre-drift followed by $+1.322$ bps of post-print move is
equally consistent with adverse selection and with the NASDAQ midpoint lagging the NBBO
while aggressive flow routes to hidden liquidity at prices that are already stale. The two
readings are observationally equivalent in single-venue data.

Re-signing and re-marking against a consolidated quote is therefore the natural next step
and, in our view, a precondition for interpreting the \emph{magnitudes} in this paper as
price discovery. We state what that work is rather than gesturing at it, since the data are
standard rather than exotic: millisecond TAQ carries the NBBO, the re-run requires
classifying the same type-5 prints against the consolidated quote in place of the NASDAQ
quote and re-marking the same horizons, and no new estimator is involved. Every magnitude
in Sections~\ref{sec:bifurcation}--\ref{sec:overstate} would be re-estimated under it. We
report the present results as provisional on that comparison.

The bifurcation is unaffected in kind: a contemporaneous comparison between two populations
observed on the same feed does not depend on whether that feed leads or lags a consolidated
one. It would, however, be \emph{sharpened} by consolidated data in one specific way. Our
at-midpoint population is defined by the NASDAQ midpoint, and a print at that midpoint may
sit away from the NBBO midpoint, where it would acquire a genuine quote-rule sign instead of
falling back on the tick rule---which \S\ref{sec:tickplacebo} shows to be a weak instrument.
Consolidated signing would therefore reclassify some of the at-midpoint population and
replace the paper's weakest classifier on its second-largest table.

\subsection{Horizon}

Every markout here is measured within a single session, at most $6.5$ hours. We therefore
describe the aggressive footprint as intraday-persistent rather than permanent, and test
separately whether it reverts over the following month.

An earlier version of this test sorted names cross-sectionally on the day's signed
aggressive-hidden flow and measured the signed return from that day's close to $k$ days
later, using non-overlapping formation dates. At $k=20$ that leaves roughly $25$ independent
cohorts, and we reported a point estimate of $-8.4$ bps with $t=-1.1$ as underpowered but
uncomfortable. That estimate turns out to be an artifact of the design rather than a weak
signal. Because the non-overlapping sample retains only every twentieth formation date, the
estimate depends on \emph{which} phase one starts from: sweeping the starting offset across
all twenty phases of the same data moves the twenty-day figure from $-36.2$ to $+25.5$ bps, a
spread of $61.6$ bps with no change to signal, weights, or estimator. The published number
was one draw from that spread.

The efficient design is a calendar-time portfolio with overlapping holds. On each day the
book holds the $k$ cohorts formed on each of the previous $k$ days at $1/k$ weight, so every
trading day contributes an observation on a diversified portfolio rather than every $k$-th
day contributing one to a sparse sort; the $k$-day cumulative return is $k$ times the mean
daily portfolio return, and inference runs on the daily series. Weights are the same
cross-sectionally demeaned, unit-gross-exposure book used throughout.

\begin{table}[htbp]
\centering
\begin{tabular}{lccc}
\toprule
\textbf{Horizon} & \textbf{Cumulative (bps)} & \textbf{NW $t$} & \textbf{Daily (bps)} \\
\midrule
1 day & $-0.69$ & $-0.55$ & $-0.688$ \\
2 days & $-0.36$ & $-0.32$ & $-0.181$ \\
3 days & $-0.50$ & $-0.37$ & $-0.166$ \\
5 days & $-1.11$ & $-0.42$ & $-0.222$ \\
10 days & $-1.08$ & $-0.34$ & $-0.108$ \\
20 days & $-1.95$ & $-0.40$ & $-0.098$ \\
\bottomrule
\end{tabular}
\caption{Multi-day behaviour after the burst day, calendar-time portfolio with overlapping
$k$-day holds (2023--2024 hidden panel, 474 names, 501 daily observations at every horizon;
Newey--West, 10 lags). No horizon shows significant reversion and every point estimate lies
within $2$ bps of zero. The corresponding non-overlapping sort at twenty days ranges from
$-36.2$ to $+25.5$ bps across starting phases of the same data; its average across phases,
$-1.1$ bps, agrees with the calendar-time estimate here.}
\label{tab:multiday}
\end{table}

Under the efficient design the multi-day picture is quiet. Every point estimate lies within
$2$ bps of zero, none approaches significance, and the twenty-day figure is $-1.95$ bps
against an intraday footprint of roughly $+2$. The reversion that the earlier design appeared
to show does not survive using all of the data rather than a twentieth of it.

We do not overstate what this settles. The standard error at twenty days is still about
$4.9$ bps, so a reversion of the order of the footprint itself remains inside the confidence
interval; what has been removed is the evidence \emph{for} such a reversion, not the
possibility of one. The honest statement is that the footprint is measured intraday, that
nothing detectable happens to it over the following month, and that the earlier suggestion
otherwise was a design artifact we have now corrected.

\subsection{Universe}

The name list is fixed by a liquidity screen among the tickers the LOBSTER archive stages,
and that staged set is itself survivor-conditioned: names that later delisted are absent
from the files we can reach rather than present as truncated series. The universe is
therefore not point-in-time.

The consequence is asymmetric across our claims and we separate them. The bifurcation of
Section~\ref{sec:bifurcation} is estimated \emph{within} name-days---aggressive against
at-midpoint prints of the same name on the same day---so a missing tail of names cannot
generate it. The cross-sectional \emph{magnitudes}, by contrast, including the level of the
footprint and the spread-scaling slope, are estimated on a panel missing exactly the
wide-spread, distressed tail where adverse selection should be most severe, and do not carry
a population interpretation. Whether the delisted names' order flow is retrievable from the
vendor is a question we have not settled.

\subsection{Attribution}

We observe anonymous message-level executions, not trader identities. There are no parent
orders, no broker records, and no linkage to holdings. We therefore make no claim about who
submits these executions, and in particular do not identify them with institutional
metaorders; where that literature is invoked it is as a comparison of magnitudes. The party
bearing the adverse-selection cost is whoever posted the hidden resting order, which may be
a designated market maker or an institution resting size, and the data cannot distinguish
them.

\subsection{Interpretation}

Finally, we separate the facts from the mechanism. That the footprint is a stable fraction
of the quoted half-spread is a measurement. That the spread is \emph{set} to compensate
providers for it is the Glosten--Milgrom interpretation, which our evidence is consistent
with but does not establish: we have not tested the causal direction between spread and
adverse selection, and both could respond to a common factor. We report the stable ratio as
a fact and the equilibrium mechanism as a reading of it.

\section{Conclusion}
\label{sec:conclusion}

Hidden liquidity on a single venue is informationally bifurcated. Executions with an
identifiable aggressor are associated with a $+2.09$ bps signed midpoint move that is flat
from three minutes to the close across the full cross-section; executions at the midpoint
carry $-0.42$ bps. Any classifier that
force-signs midpoint prints averages the two and reports something close to zero, which is
why the informativeness of hidden liquidity cannot be stated as one number. This is the
self-selection equilibrium of \citet{zhu2014} measured inside one venue and one order type,
where the confounds of a cross-venue comparison are absent.

The aggressive footprint behaves as an adverse-selection quantity. It is roughly one third
of the quoted half-spread and holds that share across a fivefold range of spread widths;
volatility, which predicts it univariately, contributes nothing once the spread is
controlled; it is $51\%$ of the effective half-spread on the Huang--Stoll accounting; it is
incremental to visible order-flow imbalance; and it survives the removal of
information-event windows. Mechanical displacement is the one alternative we cannot dismiss:
the majority of the measured level sits in the $30\%$ of prints whose touch is swept within
$100$ms and coincides with that sweep, and at that resolution displacement and
information taken through the trade are the same event.

What it is not is a tradable signal, and the reason is structural rather than incidental. A
quantity that holds a fixed proportion of the spread across the cross-section is not a
mispricing awaiting arbitrage; capturing it requires being the liquidity provider, and the
liquidity provider is already compensated by the spread that contains it. Our own attempt to
convert the footprint into a signal is consistent with this and is not reported here.

Two results limit the reading. Nearly half the signed move precedes the print on the cleanest
subset available, and a properly specified VAR conditioning on the prior return path finds no
lasting response. We therefore report an association rather than an impact. We are able to
rule out the mechanism that would tie those two facts together---decomposing the markout
against the pre-print drift leaves it essentially untouched, so it is not continuation of a
move already underway---but ruling out one alternative is not identification, and the
disagreement with the structural benchmark is left open rather than explained. And because
the book is NASDAQ's alone, we cannot separate price discovery from a local midpoint
adjusting toward a consolidated quote that moved first. The bifurcation survives all of
this; the magnitudes await a consolidated feed.

\bibliography{references}

\end{document}
